\section{\GPU{}-Native Across the Stack}
\label{sec:gpu-native}

This section breaks down what runs on the \GPU{} in \Zoo{} 4.0 and why each
component is on the device rather than the host. The design discipline is
\textsc{lp-137}'s \GPU{}-Residency Invariant: data on the hot path stays on
the device until a certificate is emitted.

\subsection{Consensus (\Quasar{} 4.0, LP-020 + LP-135)}

\Quasar{} 4.0 ships triple-consensus:

\begin{itemize}[leftmargin=1.2em]
  \item \textbf{\BLS{}12-381} --- 48-byte aggregate signature, classical
        fast path. Pairing check and miller-loop kernels run on \GPU{}.
  \item \textbf{\Corona{} (Ring-LWE)} --- post-quantum threshold proof.
        \GPU{}-resident polynomial multiplication (NTT) for signature
        share aggregation.
  \item \textbf{\MLDSA{}-65} (FIPS 204) --- post-quantum identity proof
        ($\sim$3.3 KB). Hash-and-sample, polynomial arithmetic, and
        rejection sampling all on-device.
\end{itemize}

Each layer is independently toggleable via \texttt{TripleSignRound1}; the
default 4.0 mode runs all three in parallel and emits a \Quasar{}Cert
that carries all three proofs. Photon (committee selection), Wave
(threshold voting), and Focus (confidence accumulation) execute as
device-side logic per LP-105.

\subsection{\textsc{evm} execution (LP-009)}

The \Zoo{} 4.0 \textsc{evm} is a fiber \textsc{vm}:

\begin{itemize}[leftmargin=1.2em]
  \item Each transaction maps to a \GPU{} fiber.
  \item U256 arithmetic uses 4$\times$64 limb representation in registers
        with carry-chain primitives.
  \item Cold-state \texttt{SLOAD} suspends the fiber and yields the warp
        slot to a runnable peer; the storage-fetch coalescer batches
        outstanding loads into a single device-side trie traversal.
  \item Block-STM-style optimistic concurrency
        \cite{zoo-per-llm-chains} (cf.~\textsc{lp-010}) provides parallel
        execution with deterministic re-execution on conflict.
\end{itemize}

\subsection{AI mining}

The \Zoo{} AI mining loop binds compute to receipts:

\begin{enumerate}[leftmargin=1.4em]
  \item A miner runs Zen inference inside a TEE-attested container
        (NVIDIA H100 confidential compute, AMD SEV-SNP, Intel TDX).
  \item The TEE quote plus an inference digest is submitted to the \Lux{}
        A-Chain (AIVM, \textsc{lp-134}); the A-Chain anchors the receipt.
  \item \Zoo{} reads the A-Chain receipt and attributes the inference
        to the appropriate per-\LLM{} chain (\cite{zoo-per-llm-chains}),
        updating the contributor attribution graph.
  \item Rewards are paid in \$AI to the miner and routed through the
        per-\LLM{} chain's distribution policy to upstream contributors
        (training-data, fine-tuners, evaluators).
\end{enumerate}

The hot path of this loop --- attestation verification, signature
aggregation, attribution graph update --- runs as \GPU{} kernels.

\subsection{\DEX{} matching}

\Zoo{}'s \DEX{} (see \S\ref{sec:zoo-dex}) implements \AMM{} V2/V3 natively
on the chain and exposes V4 as a precompile bridge to the \Lux{} \DEX{}
\cite{lux-dex}. \AMM{} curve math, fee calculations, tick crossings, and
order matching execute on the \GPU{}. The constant-product invariant for
V2 and the concentrated-liquidity tick math for V3 are simple enough that
a single warp can step thousands of swaps per microsecond at validator
hardware budgets.

\subsection{F-Chain \FHE{} (LP-013)}

When a \Zoo{} contract dispatches to the F-Chain \FHE{} precompile,
the ciphertext arithmetic (TFHE bootstrapping, key-switching, gadget
products) runs on the F-Chain validator's \GPU{}. Confidential ERC-20
operations (\textsc{lp-067}) --- transfer, balance check, allowance
update --- are device-side ciphertext circuits.

\subsection{Z-Chain \ZKP{} (LP-063)}

Verification of Groth16, Halo2, and Plonky2 proofs runs as \GPU{}
kernels on the Z-Chain validator. \Zoo{} dispatches verification
requests via precompile; the Z-Chain returns a verification receipt
that the calling \Zoo{} contract treats as a boolean.

\subsection{The discipline}

The common pattern across all of the above: data lands on the \GPU{}
when a transaction enters the mempool, executes its lifecycle on
the device (verify, run, prove, log), and only leaves the device
when a certificate is committed to disk or sent over the wire.
The CPU becomes the orchestration layer; the \GPU{} is where the
ledger lives during the block.
